% ---------- Язык и кодировка (pdfLaTeX) ----------
\usepackage{cmap}              % копируемый/ищущийся русский текст в PDF (до fontenc)
\usepackage[T2A]{fontenc}
\usepackage[utf8]{inputenc}
\usepackage[english,russian]{babel}

% ---------- Вёрстка ----------
\usepackage[a4paper, margin=2.5cm]{geometry}
\usepackage{microtype}
\usepackage{enumitem}
\setlist{nosep}
\usepackage{booktabs, longtable, array}
\usepackage{indentfirst}       % абзацный отступ после заголовка, как принято в русском

% ---------- Математика ----------
\usepackage{amsmath, amssymb, amsthm}
\usepackage{thmtools}          % \declaretheorem: корректные имена в \cref при общем счётчике
\usepackage{mathtools}
\usepackage{bm}

% ---------- Графика ----------
\usepackage{graphicx}
\graphicspath{{figures/}}
\usepackage{tikz}
\usepackage{pgfplots}
\pgfplotsset{compat=1.18}

% ---------- Цвета и рамки ----------
\usepackage{xcolor}
\usepackage[most]{tcolorbox}

% ---------- Ссылки (подключать последними) ----------
\usepackage{hyperref}
\hypersetup{
  colorlinks = true,
  linkcolor  = blue!60!black,
  urlcolor   = blue!60!black,
  citecolor  = green!40!black,
  pdftitle   = {Математическая статистика. Конспект лекций},
}
\usepackage[nameinlink]{cleveref}
